\documentclass[11pt]{article}
\usepackage[margin=1in]{geometry}
\usepackage{xcolor}
\usepackage{tcolorbox}
\usepackage{hyperref}
\usepackage{enumitem}
\usepackage{booktabs}
\usepackage{graphicx}
\usepackage{longtable}

\definecolor{osaorange}{RGB}{220,115,38}
\definecolor{aiblue}{RGB}{30,144,255}

\hypersetup{
    colorlinks=true,
    linkcolor=blue,
    urlcolor=blue,
    citecolor=blue
}

\title{\textbf{H524 Introduction to Biostatistics\\
\Large Final Project: Learn and Apply Statistics with AI\\
Fall 2025}}
\author{Dr. John Molitor}
\date{}

\begin{document}
\maketitle

\begin{tcolorbox}[colback=aiblue!10,colframe=aiblue,boxrule=2pt,arc=2pt,left=6pt,right=6pt,top=6pt,bottom=6pt]
\centering
\textbf{\large Project Overview}\\[0.3em]
Use AI to teach yourself one statistical method we learned this semester — and apply it to a real question. You'll then teach the class what you learned.
\end{tcolorbox}

\section{Your Task}

Teams of 4-5 students will:

\begin{enumerate}
\item Choose a core stats concept (e.g., chi-square test, correlation, regression, t-test, ANOVA, bootstrap, etc.)
\item Learn and apply it using AI tools
\item Compare at least two AI platforms (demo both chat and code)
\item Show what you learned, what worked, what didn't, and how you verified the AI's answers
\end{enumerate}

\section{Teams}

\textbf{Size:} 4-5 students per team

\textbf{Formation:} Self-select teams by Week 4. Post team composition on Canvas discussion board.

\textbf{Roles:} While all team members contribute to all aspects, designate:
\begin{itemize}
\item \textbf{Project Coordinator:} Manage timeline and submissions
\item \textbf{AI Tool Lead:} Coordinate AI platform testing and documentation
\item \textbf{Data/Code Lead:} Ensure reproducibility and code quality
\item \textbf{Presentation Lead:} Coordinate slide development and presentation flow
\item \textbf{Documentation Lead:} Compile prompt logs and write-up
\end{itemize}

\section{Statistical Concepts to Choose From}

Select ONE of the following core methods from H524:

\subsection{Option 1: Chi-Square Test}
Test for independence or goodness-of-fit with categorical data.

\textbf{Example Topics:}
\begin{itemize}
\item Is music genre independent of mood tag?
\item Do voting patterns differ by age group?
\item Is disease prevalence independent of geographic region?
\end{itemize}

\subsection{Option 2: Correlation Analysis}
Examine linear relationships between continuous variables.

\textbf{Example Topics:}
\begin{itemize}
\item Relationship between study hours and exam scores
\item Air quality index vs respiratory hospital admissions
\item Social media usage and self-reported well-being
\end{itemize}

\subsection{Option 3: Linear Regression}
Model continuous outcomes with one or more predictors.

\textbf{Example Topics:}
\begin{itemize}
\item Do housing prices depend on square footage and location?
\item Predicting blood pressure from age and BMI
\item Effect of exercise duration on weight loss
\end{itemize}

\subsection{Option 4: t-Test (Independent or Paired)}
Compare means between two groups.

\textbf{Example Topics:}
\begin{itemize}
\item Do organic apples cost more than conventional ones?
\item Sleep quality before vs after intervention
\item Math scores: traditional vs flipped classroom
\end{itemize}

\subsection{Option 5: ANOVA}
Compare means across three or more groups.

\textbf{Example Topics:}
\begin{itemize}
\item Pain scores across three medication types
\item Plant growth under different fertilizer treatments
\item Test anxiety levels across different study methods
\end{itemize}

\subsection{Option 6: Bootstrap Confidence Intervals}
Use resampling to estimate parameter uncertainty.

\textbf{Example Topics:}
\begin{itemize}
\item What's a 95\% CI for average delivery time?
\item Bootstrap CI for correlation coefficient
\item Median household income CI for small samples
\end{itemize}

\section{AI Tool Requirements}

\textbf{You must compare at least 2 of these:}
\begin{itemize}
\item ChatGPT (with or without code interpreter)
\item Anthropic Claude
\item Google Gemini
\item Microsoft Copilot
\item Other AI tools (with instructor approval)
\end{itemize}

\textbf{What to demonstrate:}
\begin{itemize}
\item Show prompts you used with each AI
\item Compare how each AI explained or solved your problem
\item Explain what you trusted, corrected, or changed
\item Demonstrate both chat-based help AND code generation (where applicable)
\end{itemize}

\section{Timeline}

\begin{center}
\begin{tabular}{ll}
\toprule
\textbf{Week} & \textbf{Milestone} \\
\midrule
Week 4 & Teams formed, posted on Canvas \\
Week 5 & \textbf{150-word proposal due Sunday 11:59 PM} \\
Week 6-7 & In-class proposal discussions and feedback \\
Week 8-9 & Work on project (office hours available) \\
Week 10 & \textbf{Presentations + all materials due} \\
\bottomrule
\end{tabular}
\end{center}

\section{What to Turn In}

\subsection{1. Presentation (10 minutes total)}

Use slides to cover:

\textbf{Your Question} (1-2 min)
\begin{itemize}
\item What did you want to find out?
\item Why is it interesting or relevant?
\end{itemize}

\textbf{Your Data} (1 min)
\begin{itemize}
\item Where it came from and what's in it
\item Sample size, variables, any preprocessing
\end{itemize}

\textbf{Your Method} (2 min)
\begin{itemize}
\item Which stats concept you chose and why
\item Brief explanation of the method
\end{itemize}

\textbf{Your AI Process} (3-4 min)
\begin{itemize}
\item Show the prompts or examples you tried
\item Compare how at least 2 AIs explained or solved it
\item Explain what you trusted, corrected, or changed
\item Demo both chat assistance and code generation
\end{itemize}

\textbf{Your Results} (2 min)
\begin{itemize}
\item Key outputs, graphs, and interpretation
\item Statistical significance and practical meaning
\end{itemize}

\textbf{Takeaways} (1 min)
\begin{itemize}
\item What classmates should learn from your process
\item Recommendations for using AI in statistics
\end{itemize}

\subsection{2. Short Write-Up (2 pages max)}

Include:
\begin{itemize}
\item Which AI tools you used and why
\item What each AI did well/poorly
\item One workflow or strategy you'd recommend for others
\item Any code, screenshots, or outputs that show your steps
\end{itemize}

\subsection{3. Prompt Log + Data/Code}

Attach:
\begin{itemize}
\item Prompts and short notes (copy/paste or screenshots)
\item Dataset (or link to data source)
\item Code or notebook (if applicable) with clear comments
\item 2-3 sentences on each teammate's role
\end{itemize}

\section{Grading (100 points total)}

\begin{center}
\begin{tabular}{lcp{8cm}}
\toprule
\textbf{Category} & \textbf{Points} & \textbf{Description} \\
\midrule
Question \& Data & 10 & Clear, relevant, ethical use of data \\
Method \& Accuracy & 25 & Correct test/model, assumptions checked \\
AI Comparison & 25 & 2+ AIs, transparent logs, clear reflection \\
Interpretation & 15 & Results explained clearly and correctly \\
Reproducibility & 15 & Code/data provided, process repeatable \\
Presentation & 10 & Organized, engaging, within time \\
\midrule
\textbf{Total} & \textbf{100} & \\
\bottomrule
\end{tabular}
\end{center}

\textbf{Deductions:}
\begin{itemize}
\item -5 pts for going over time or missing files
\item -20 pts if fewer than 2 AIs are compared
\end{itemize}

\section{Presentation Signup}

30 students $\rightarrow$ roughly 6-8 groups $\rightarrow$ presented over 2 class sessions in Week 10

Each slot = 10 min presentation + 2 min Q\&A

Sign up on the shared sheet (first come, first served). Swaps allowed with instructor approval.

\begin{center}
\begin{tabular}{lll}
\toprule
\textbf{Session} & \textbf{Time} & \textbf{Group/Topic} \\
\midrule
Session A & 00:00-00:12 & \\
Session A & 00:12-00:24 & \\
Session A & 00:24-00:36 & \\
Session A & 00:36-00:48 & \\
Session A & 00:48-01:00 & \\
Session A & 01:00-01:12 & \\
\midrule
Session B & 00:00-00:12 & \\
Session B & 00:12-00:24 & \\
Session B & 00:24-00:36 & \\
Session B & 00:36-00:48 & \\
Session B & 00:48-01:00 & \\
Session B & 01:00-01:12 & \\
\bottomrule
\end{tabular}
\end{center}

See separate signup sheet for detailed scheduling.

\section{Tips for Success}

\begin{enumerate}
\item \textbf{Start Early:} Don't wait until Week 9 to begin
\item \textbf{Document as You Go:} Save all prompts and AI responses immediately
\item \textbf{Verify Everything:} Never trust AI output without checking
\item \textbf{Practice Your Presentation:} Time it and rehearse transitions
\item \textbf{Use Simple Data:} Don't choose overly complex datasets
\item \textbf{Ask for Help:} Use office hours if you're stuck
\item \textbf{Be Honest:} Show what didn't work, not just successes
\end{enumerate}

\section{Example Project Flow}

\subsection{Week 5: Proposal Example}
\textit{"Our team will investigate whether there is a relationship between daily screen time and sleep quality among college students using correlation analysis. We will collect or use existing survey data with variables for hours of screen time and self-reported sleep scores. We plan to use ChatGPT and Claude to help us understand correlation assumptions, generate R code, and interpret results. Team members: Alex (coordinator), Jordan (AI tools), Sam (data/code), Taylor (presentation), Morgan (documentation)."}

\subsection{Week 6-7: What You're Doing}
\begin{itemize}
\item Finding or collecting your data
\item Testing different AI prompts
\item Learning about your statistical method from multiple sources
\item Starting to code and analyze
\end{itemize}

\subsection{Week 10: What You're Showing}
\begin{itemize}
\item Side-by-side comparison of AI responses
\item Your final analysis with interpretation
\item Lessons learned about AI strengths and weaknesses
\item Recommendations for your classmates
\end{itemize}

\section{Academic Integrity}

\begin{tcolorbox}[colback=red!5,colframe=red,boxrule=1pt,arc=2pt,left=6pt,right=6pt,top=6pt,bottom=6pt]
\textbf{Important:}
\begin{itemize}
\item Using AI is required and encouraged for this project
\item You must be transparent about what AI generated vs what you created
\item All AI usage must be documented in your prompt log
\item You are responsible for verifying all AI outputs for accuracy
\item Each team member must contribute meaningfully
\item Data fabrication or plagiarism will result in project failure
\end{itemize}
\end{tcolorbox}

\section{Resources}

\subsection{Getting Started with AI Tools}
\begin{itemize}
\item ChatGPT: \url{https://chat.openai.com}
\item Claude: \url{https://claude.ai}
\item Gemini: \url{https://gemini.google.com}
\item Microsoft Copilot: \url{https://copilot.microsoft.com}
\end{itemize}

\subsection{Finding Data}
See \texttt{dataset\_suggestions.md} for curated list of accessible datasets.

\subsection{R Resources}
\begin{itemize}
\item Course materials on Canvas
\item R Documentation: \url{https://www.rdocumentation.org}
\item Quick-R: \url{https://www.statmethods.net}
\end{itemize}

\subsection{Statistical Concepts Review}
\begin{itemize}
\item Course textbook chapters
\item Khan Academy Statistics
\item Course lecture slides and recordings
\end{itemize}

\section{Frequently Asked Questions}

\textbf{Q: Can we use a statistical method not listed?}\\
A: Check with the instructor first. It should be a core H524 concept.

\textbf{Q: What if the AIs give different answers?}\\
A: Perfect! That's exactly what you should discuss in your presentation.

\textbf{Q: How technical should the write-up be?}\\
A: Focus on the AI comparison and learning process, not just the statistics.

\textbf{Q: Can we use our own data?}\\
A: Yes, as long as it's appropriate and you have permission to use it.

\textbf{Q: What if we can't get the analysis to work?}\\
A: Document what went wrong and what you learned. Troubleshooting is valuable.

\textbf{Q: Do we all present together?}\\
A: Yes, but divide the presentation so everyone speaks roughly equally.

\end{document}

